\section{Síntesis y ampliación: cuatro desigualdades y sus condiciones de aplicación}

Al final, la conferencia condensa su línea principal en cuatro comparaciones concisas: reasoning es mejor que no reasoning, RL finetuning es mejor que SFT, agregar varias respuestas es mejor que una sola respuesta, y retrieval + reasoning es mejor que reasoning por sí solo. Deben leerse como \textbf{direcciones de diseño dentro de las tareas y los presupuestos que trató la conferencia}, no como teoremas matemáticos incondicionales. Cada «mejor» se obtiene a cambio de costo, infraestructura y nuevos modos de falla.

\subsection{Cómo se traducen las cuatro conclusiones de la clase en la pila del sistema}

La slide 47 es una página de resumen muy compacta; esta sección la despliega en una pila de sistema implementable. Al leer la figura, observe primero que los cuatro «signos de mayor que» modifican, respectivamente, la longitud del razonamiento, la distribución de los parámetros, la estadística del muestreo y la información externa; después revise qué costo y qué responsabilidad de verificación añade cada capa. Las cuatro intervenciones pueden combinarse en un mismo sistema; no se trata de elegir una de cuatro.

\lecturefigure{slide-47.jpg}{Resumen de la clase: reasoning, RL finetuning, aggregation y retrieval}{47}

\begin{center}
\begin{tabularx}{0.97\textwidth}{p{0.18\textwidth}p{0.23\textwidth}X X}
\toprule
Capa de intervención & Qué cambia & Beneficio principal & Riesgo nuevo \\
\midrule
Reasoning tokens & Cómputo secuencial de una sola generación & Descomposición, estado temporal, pasos verificables & Latencia, redundancia, acumulación de errores \\
RL finetuning & Distribución base de trayectorias y respuestas & Las rutas exitosas aparecen con más frecuencia & Sesgo del verifier, reward hacking \\
Aggregation & Estimación de la respuesta tras varios muestreos & Reduce el azar de una sola ruta & Costo en tokens, errores correlacionados, análisis sintáctico \\
Retrieval & Evidencia externa visible para el modelo & Conocimiento actualizado, casos, principios & Ruido, obsolescencia, riesgos de procedencia y de inyección \\
\bottomrule
\end{tabularx}
\end{center}

\begin{importantbox}{Una sola expresión que resume todo el sistema}
Un reasoning system práctico puede escribirse así: primero se recupera $z$, luego se muestrean de la política varias trayectorias $r^{(1:K)}$ y respuestas $a^{(1:K)}$, y se usan el verifier y la aggregation para elegir la salida final. La capacidad proviene de la combinación de componentes; la confiabilidad depende de que la interfaz más débil sea observable, verificable y reproducible.
\end{importantbox}

\subsection{El siguiente avance: más allá de la respuesta única verificable}

La slide 48 dirige los problemas abiertos hacia dos frentes: resolver tareas que no tienen una respuesta única verificable, y construir aplicaciones reales en lugar de seguir saturando benchmarks. Lo primero exige pasar de la binary correctness a la evaluación multiobjetivo, la evidencia del proceso, la retroalimentación a largo plazo y la expresión de la incertidumbre; lo segundo exige registrar de extremo a extremo el valor para el usuario, la gravedad de las fallas, el costo, la latencia, los permisos y la capacidad de recuperación.

\lecturefigure{slide-48.jpg}{La siguiente etapa: tareas sin respuesta única verificable y aplicaciones reales}{48}

\teachervoice{00:56:31--00:57:30, el docente dice que los benchmarks se saturan muy rápido y que prefiere ver aplicaciones reales y tareas que vayan más allá de la respuesta única verificable. Nota de clase: esto no niega los benchmarks, sino que exige que funcionen como herramientas de diagnóstico y no como sustitutos del valor final del producto.}

Las tareas abiertas pueden adoptar una aceptación en varias capas, en lugar de obligar a un solo reward a cargar con todo el valor:
\begin{enumerate}
  \item \textbf{Restricciones duras}: fuentes de los hechos, permisos, formato, seguridad y límites legales;
  \item \textbf{Métricas de la tarea}: completitud, ejecutabilidad, objetivos del usuario y costo;
  \item \textbf{Evidencia del proceso}: fuentes recuperadas, trace de herramientas, contraejemplos e incertidumbre;
  \item \textbf{Resultados a largo plazo}: adopción real, tasa de corrección de errores, reversiones y externalidades negativas.
\end{enumerate}

\subsection{Q\&A: seis salvedades que no pueden quitarse de las conclusiones de la clase}

Los nueve minutos de preguntas y respuestas posteriores a la conferencia añadieron las incertidumbres más importantes que faltaban en las slides. Primero, en la experiencia del ponente, el answer-token log-prob puede servir como señal de hallucination, pero no se presentó como un detector perfecto. Segundo, search puede usarse como herramienta que el modelo invoca; excluir el search externo al «estudiar reasoning» no significa que el producto no deba usar search. Tercero, todavía falta una explicación satisfactoria de cómo el entrenamiento remodela la distribución de la secuencia completa, y la intuición sobre el top-token no puede extenderse hasta convertirse en una teoría completa.

\teachervoice{00:57:30--01:01:00, el docente responde sobre confidence y search: el log-prob es útil en la práctica, pero search es una herramienta integrable; a él le interesa lo que el modelo mismo puede aprender sin depender de una enumeración exhaustiva explícita. Advertencia de los apuntes: aislar variables en la investigación y buscar el mejor desempeño en un producto son dos objetivos distintos.}

Cuarto, la frontera entre reasoning y answer requiere un parser; si la respuesta final es un programa, la extracción y la verificación se vuelven claramente más difíciles. Quinto, un candidato de baja confianza a veces puede ser correcto, y self-consistency también puede fallar. Sexto, ante la pregunta «AGI en dos a cinco años: ¿qué deberían aprender los niños?», el docente se negó a aceptar la premisa del calendario y subrayó que los métodos actuales aún tienen muchas limitaciones y que lo que realmente falta son killer applications; reconoció expresamente el valor de los asistentes de programación, pero no equiparó los productos de chat con una AGI ya llegada.

\teachervoice{01:01:00--01:06:04, el docente admite que la distribution reshaping todavía no se entiende bien, que las respuestas en forma de programa requieren un análisis sintáctico cuidadoso y que self-consistency no es perfecta, y se mantiene escéptico ante los calendarios de AGI a corto plazo. El criterio práctico final es si el modelo da lugar a aplicaciones realmente útiles, no cuánta atención recibe la discusión.}

\begin{warningbox}{Límites de la evidencia de la clase de 2025}
En esta conferencia, Gemini 2.0 thinking, la aggregation de o1, la UI de Deep Research, las cifras de GSM8K y los juicios sobre productos deben entenderse en el momento de la clase, abril de 2025. Los apuntes conservan los mecanismos y las lecciones de ingeniería, sin reescribir las capturas de pantalla y las afirmaciones empíricas de entonces como especificaciones de productos actuales de 2026 ni como garantías generales.
\end{warningbox}

\subsection{Doce preguntas antes de iniciar un proyecto de reasoning}

\small
\begin{multicols}{2}
\begin{enumerate}
  \item ¿Qué estados intermedios designa concretamente «Reasoning» en esta tarea?
  \item Si la respuesta directa falla, ¿falta capacidad o falla el ordenamiento del decoding?
  \item ¿Puede analizarse de forma estable la frontera de la respuesta, y cómo se normalizan las respuestas equivalentes?
  \item ¿Las trayectorias humanas cubren los estados que el modelo realmente visitará?
  \item ¿Qué brecha de aproximación hay entre el reward y la calidad real?
  \item ¿Qué tan grandes son los false positive y false negative del verifier?
  \item ¿Qué se escala: la longitud de la trayectoria, el número de muestras, el árbol de búsqueda o las llamadas a herramientas?
  \item ¿Las múltiples muestras son suficientemente independientes, o comparten el mismo sesgo sistemático?
  \item ¿La consistency está calibrada sobre la distribución objetivo?
  \item ¿Cómo se conservan la procedencia, la fecha y los conflictos de lo recuperado con retrieval?
  \item ¿Cómo combinan las tareas abiertas restricciones duras, preferencias y retroalimentación a largo plazo?
  \item Tras el despliegue, ¿cómo se reproduce la trace, se detectan regresiones y se detiene o revierte el sistema?
\end{enumerate}
\end{multicols}
\normalsize

\subsection{Lecturas adicionales}

\scriptsize
\begin{multicols}{2}
\noindent\textbf{Cómputo secuencial: }Li et al., \emph{Chain of Thought Empowers Transformers to Solve Inherently Serial Problems}. \href{https://arxiv.org/abs/2402.12875}{2402.12875}.\par
\noindent\textbf{Decodificación sin prompt: }Wang and Zhou, \emph{Chain-of-Thought Reasoning Without Prompting}. \href{https://arxiv.org/abs/2402.10200}{2402.10200}.\par
\noindent\textbf{Agregación de respuestas: }Wang et al., \emph{Self-Consistency Improves Chain of Thought Reasoning in Language Models}. \href{https://arxiv.org/abs/2203.11171}{2203.11171}.\par
\noindent\textbf{Agregación de texto libre: }Chen et al., \emph{Universal Self-Consistency for Large Language Model Generation}. \href{https://arxiv.org/abs/2311.17311}{2311.17311}.\par
\end{multicols}
\normalsize

\end{document}
